\documentclass[11pt]{article}
\usepackage{amsmath,amssymb,amsthm}
\usepackage{hyperref}
\usepackage[margin=1in]{geometry}

\newcommand{\N}{\mathbb{N}}
\newcommand{\Z}{\mathbb{Z}}
\newcommand{\Q}{\mathbb{Q}}
\newcommand{\La}{\Lambda}
\newcommand{\GK}{\operatorname{GK}}
\newcommand{\CPP}{\operatorname{CPP}}
\newcommand{\ALCD}{\operatorname{ALCD}}
\newcommand{\IP}{\operatorname{IP}}
\newcommand{\SSYT}{\operatorname{SSYT}}
\newcommand{\qpoch}[2]{(#1;#2)_\infty}
\newcommand{\qpochn}[3]{(#1;#2)_{#3}}
\newcommand{\qbin}[2]{\genfrac{[}{]}{0pt}{}{#1}{#2}}
\DeclareMathOperator{\tr}{tr}
\DeclareMathOperator{\wt}{wt}
\DeclareMathOperator{\sgn}{sgn}
\DeclareMathOperator{\Fix}{Fix}

\theoremstyle{plain}
\newtheorem{theorem}{Theorem}
\newtheorem{conjecture}[theorem]{Conjecture}
\newtheorem{proposition}[theorem]{Proposition}
\newtheorem{lemma}[theorem]{Lemma}
\newtheorem{corollary}[theorem]{Corollary}
\theoremstyle{definition}
\newtheorem{definition}[theorem]{Definition}
\newtheorem{example}[theorem]{Example}
\newtheorem{remark}[theorem]{Remark}

\title{$Q$-Local: Bender--Knuth Involutions \\
  and $Q$-Row Operations}
\author{Robin Langer}
\date{\today}

\begin{document}
\maketitle

\begin{abstract}
The classical Sch\"utzenberger involution decomposes into
Bender--Knuth involutions $\sigma_i$, each of which swaps two
adjacent rows of the growth diagram.  We give a self-contained
account of these involutions, their connection to local rules,
and their $Q$-analogues.  The $Q$-Bender--Knuth involution
$\sigma_i^Q$ acts on ground-set elements
$(S, \mathfrak{d}) \in \mathcal{G}_{n,c}$ by switching the
cascade direction at the $i$-th diagonal of the ALCD\@.  The
condition $(\sigma_i^Q)^2 = \mathrm{id}$ is the atomic
requirement for the $Q$-Sch\"utzenberger
identity~\cite{Q_SCHUTZ}.
\end{abstract}

\tableofcontents


%% ============================================================
\part{The Classical Theory}
%% ============================================================


%% ============================================================
\section{The Bender--Knuth involutions}
\label{sec:BK-def}
%% ============================================================

Let $T$ be a semistandard Young tableau (SSYT) of shape
$\lambda$ with entries in $\{1, 2, 3, \ldots\}$.

\begin{definition}[Free and paired entries]
\label{def:free-paired}
Fix an index $i \geq 1$.  Consider the entries $i$ and
$i + 1$ in the tableau~$T$.  In each row~$r$, suppose there
are $a_r$ copies of~$i$ and $b_r$ copies of~$i + 1$.  The
column-strictness condition forces a certain number of these
to be \emph{paired}: an $i$ in column~$c$ is paired with the
$i+1$ in column~$c$ of the row above (if it exists), and
conversely.  More precisely:
\begin{itemize}
  \item The entries $i$ and $i+1$ in row~$r$ occupy a
    contiguous block of columns (by the semistandard
    condition: all $i$'s come before all $(i+1)$'s).
  \item An $i$ is \emph{paired} if the cell directly below it
    (in row $r+1$) contains $i+1$.
  \item An $i+1$ is \emph{paired} if the cell directly above
    it (in row $r-1$) contains~$i$.
  \item All other $i$'s and $(i+1)$'s are \emph{free}.
\end{itemize}
Let $f_r$ be the number of free $i$'s and $g_r$ the number of
free $(i+1)$'s in row~$r$.
\end{definition}

\begin{definition}[Bender--Knuth involution $\sigma_i$]
\label{def:BK}
The involution $\sigma_i \colon \SSYT(\lambda) \to
\SSYT(\lambda)$ acts by:
\begin{enumerate}
  \item Keep all paired entries fixed.
  \item In each row~$r$, exchange the free entries:
    replace each free~$i$ by~$i + 1$ and each free~$i+1$
    by~$i$.
\end{enumerate}
Equivalently, in each row, swap $(f_r, g_r) \mapsto
(g_r, f_r)$.
\end{definition}

\begin{proposition}[Properties of $\sigma_i$]
\label{prop:BK-properties}
\mbox{}
\begin{enumerate}
  \item $\sigma_i^2 = \mathrm{id}$ \textup{(involution)}.
  \item $\sigma_i$ preserves the shape~$\lambda$.
  \item $\sigma_i$ exchanges the multiplicities of~$i$
    and~$i+1$ in the content.
  \item $\sigma_i \circ \sigma_j = \sigma_j \circ \sigma_i$
    for $|i - j| \geq 2$
    \textup{(distant involutions commute)}.
\end{enumerate}
\end{proposition}

\begin{proof}
(1)~Swapping $(f_r, g_r) \mapsto (g_r, f_r)$ twice gives the
identity.
(2)~Only entries $i \leftrightarrow i+1$ are changed; the
shape is determined by row lengths, which are unaffected.
(3)~Free $i$'s become free $(i+1)$'s and vice versa; paired
entries are unchanged.
(4)~The entries $\{i, i+1\}$ and $\{j, j+1\}$ are disjoint
when $|i-j| \geq 2$, so $\sigma_i$ and $\sigma_j$ act on
independent sets of cells.
\end{proof}

\begin{corollary}[Schur functions are symmetric]
\label{cor:schur-symmetric}
The Schur function
$s_\lambda = \sum_{T \in \SSYT(\lambda)}
x^{\mathrm{content}(T)}$ is symmetric: $\sigma_i$ gives a
weight-preserving, content-swapping bijection on
$\SSYT(\lambda)$, so the coefficient of $x^\mu$ equals the
coefficient of~$x^{s_i(\mu)}$.
\end{corollary}


%% ============================================================
\section{Connection to growth diagrams and local rules}
\label{sec:growth}
%% ============================================================

\subsection{The growth diagram}

A \emph{growth diagram} of size $m \times n$ is an array of
partitions $\lambda^{(i,j)}$ ($0 \leq i \leq m$,
$0 \leq j \leq n$) satisfying:
\begin{itemize}
  \item Boundary conditions on $\lambda^{(i,0)}$ and
    $\lambda^{(0,j)}$.
  \item A \emph{local rule} at each face $(i,j)$: the diamond
    $(\lambda^{(i-1,j-1)};\; \lambda^{(i-1,j)},\,
    \lambda^{(i,j-1)};\; \lambda^{(i,j)})$ with face label
    $k_{i,j} \geq 0$.
\end{itemize}
Each diagonal~$i$ of the growth diagram corresponds to one
application of
$T_n P_m = \sum_k P_{m-k} T_{n-k}$.
The face labels along diagonal~$i$ form the $i$-th row of the
ALCD\@.


\subsection{Row-swapping $=$ Bender--Knuth}

The key insight (van~Leeuwen~\cite{vanLeeuwen}):

\begin{proposition}
\label{prop:row-swap}
Swapping diagonals $i$ and $i+1$ of the growth
diagram---i.e., recomputing the vertex partitions with the
data from diagonals $i$ and $i+1$ interchanged---is equivalent
to applying the Bender--Knuth involution $\sigma_i$ to the
corresponding tableau.
\end{proposition}

The dictionary:
\begin{center}
\begin{tabular}{l|l}
\textbf{Tableaux} & \textbf{Growth diagrams} \\ \hline
$\sigma_i$: swap entries $i \leftrightarrow i+1$ &
  swap diagonals $i \leftrightarrow i+1$ \\
free entries in a row &
  face labels that change under the swap \\
paired entries &
  face labels that stay fixed \\
$\sigma_i^2 = \mathrm{id}$ &
  double swap $=$ identity
\end{tabular}
\end{center}


\subsection{Burge vs.\ RSK at one diagonal}

The two local rules (Burge and RSK) differ only in cascade
direction:
\begin{itemize}
  \item \textbf{Burge forward}: cascade \textbf{up}
    ($\varepsilon(i) = \min\{j > i : j \notin
    A \cup B \cup C\}$).
  \item \textbf{RSK forward}: cascade \textbf{down}
    ($\delta(i) = \max\{j < i : j \notin
    A \cup B \cup C\}$).
\end{itemize}

\begin{proposition}
\label{prop:switch-one-diagonal}
Building a growth diagram using Burge at all diagonals, then
switching to RSK at diagonal~$i$ alone, produces a growth
diagram that differs from the original by $\sigma_i$ applied
to the face-label data.
\end{proposition}


\subsection{Sch\"utzenberger $=$ product of Bender--Knuth}

The Sch\"utzenberger involution
$\psi = \psi_R \circ \psi_B^{-1}$ switches \emph{every}
diagonal from Burge to RSK\@.  By
Proposition~\ref{prop:switch-one-diagonal}:
\begin{equation}\label{eq:psi-decomp}
  \psi = \prod_{i} \sigma_i,
\end{equation}
where the product is over all diagonals.

The classical proof that $\psi^2 = \mathrm{id}$ uses:
\begin{enumerate}
  \item Each $\sigma_i^2 = \mathrm{id}$.
  \item The specific order of the product and the growth
    diagram structure ensure the double application is
    the identity.
\end{enumerate}

\begin{remark}[Subtlety]
\label{rem:subtlety}
The identity $\psi^2 = \mathrm{id}$ does \emph{not} follow
formally from $\sigma_i^2 = \mathrm{id}$ alone, because
adjacent $\sigma_i$ and $\sigma_{i+1}$ do not commute in
general.  The deeper reason is that $\psi$ is a global
Burge$\leftrightarrow$RSK toggle: switching all diagonals
twice restores the original rule at every diagonal.
\end{remark}


%% ============================================================
\part{The $Q$-Enriched Theory}
%% ============================================================


%% ============================================================
\section{$Q$-row operations on the ground set}
\label{sec:Q-row}
%% ============================================================

\subsection{The ALCD as a sequence of diagonals}

An ALCD $\mathfrak{d} \in \ALCD_n(c)$ encodes a growth
diagram of $n$~diagonals.  We write the face labels along
diagonal~$i$ as $\mathfrak{d}_i = (k_{i,1}, k_{i,2}, \ldots)$
and the total weight as
$|\mathfrak{d}| = \sum_{i,j} k_{i,j}$.

Each diagonal~$i$ corresponds to one commutation step.
Switching the rule (Burge $\leftrightarrow$ RSK) at
diagonal~$i$ changes the vertex partitions at that diagonal
and, in the $Q$-enriched case, propagates to all subsequent
diagonals.


\subsection{The $Q$-row switch}

\begin{definition}[$Q$-row switch $\sigma_i^Q$]
\label{def:Q-row-switch}
Let $(S, \mathfrak{d}) \in \mathcal{G}_{n,c}$ with
$\mathfrak{d} \in \ALCD_{n-|S|}(c)$.  The
\emph{$Q$-row switch at diagonal~$i$}
($1 \leq i \leq n - |S|$) is:
\begin{enumerate}
  \item Reconstruct the growth diagram for $\mathfrak{d}$ up
    to diagonal~$i-1$ using the current rule (say Burge).
  \item At diagonal~$i$, \textbf{switch} the rule: recompute
    the vertex partitions and face labels along diagonal~$i$
    using the opposite cascade direction.
  \item \textbf{Recompute} all subsequent diagonals
    $i+1, i+2, \ldots$ using the \emph{original} rule with
    the modified vertex partitions from step~(2).
\end{enumerate}
The result is $\sigma_i^Q(S, \mathfrak{d})
= (S,\, \mathfrak{d}')$ where
$\mathfrak{d}' \in \ALCD_{n-|S|}(c)$.
\end{definition}


%% ============================================================
\section{The propagation phenomenon}
\label{sec:propagation}
%% ============================================================

\subsection{Classical: strictly local}

In the standard growth diagram, switching the rule at
diagonal~$i$ changes the vertex partitions at diagonal~$i$
but does \textbf{not} change the face labels at any other
diagonal.  This is because:
\begin{itemize}
  \item The face-label generating function is $1/(1-t)$
    everywhere (universal, independent of strip sizes).
  \item The face labels at diagonal~$i' > i$ depend only on
    the vertex partitions at diagonal~$i'-1$, which are
    recomputed using the original rule and happen to give the
    same face labels (because $1/(1-t)$ doesn't depend on
    which partitions appear).
\end{itemize}

\subsection{$Q$-enriched: propagation through strip sizes}

In the $Q$-enriched growth diagram, the face label~$k$ at a
face with strip sizes $(a,b)$ carries weight
$Q_{a+k} \cdot Q_{b+k}$.  The strip sizes depend on the
vertex partitions, which change when we switch the rule at
diagonal~$i$.

\begin{proposition}[Propagation]
\label{prop:propagation}
Switching the rule at diagonal~$i$ in the $Q$-enriched growth
diagram:
\begin{enumerate}
  \item Changes the vertex partitions at diagonal~$i$
    (same as classical).
  \item Changes the strip sizes $(a', b')$ at diagonals
    $i+1, i+2, \ldots$ (because they depend on the modified
    vertex partitions).
  \item Changes the $Q$-weights at diagonals $> i$ (because
    $Q_{a'+k'} \cdot Q_{b'+k'}$ depends on the new strip
    sizes).
  \item May change the face labels $k_{i',j}$ at diagonals
    $> i$ (because the local rule with different inputs
    produces different outputs).
\end{enumerate}
\end{proposition}

This is the \emph{propagation phenomenon}: a local switch at
one diagonal cascades through the entire subsequent portion
of the growth diagram.  It is the same obstruction
identified in~\cite{LOCAL_RULE}, Sections~4--6, for the full
Sch\"utzenberger involution, now localised to a single
diagonal switch.


%% ============================================================
\section{$(\sigma_i^Q)^2 = \mathrm{id}$: proof}
\label{sec:involution-proof}
%% ============================================================

Despite the propagation, the $Q$-row switch is an involution.

\begin{theorem}
\label{thm:Q-sigma-involution}
$(\sigma_i^Q)^2 = \mathrm{id}$ on $\mathcal{G}_{n,c}$.
\end{theorem}

\begin{proof}
Let $(S, \mathfrak{d})$ be a ground-set element.  Apply
$\sigma_i^Q$:
\begin{enumerate}
  \item[(a)] Switch diagonal~$i$ from Burge to RSK\@.  The
    vertex partitions at diagonal~$i$ change.  Call the
    modified vertex partitions~$\mu^{(i,\cdot)}$.
  \item[(b)] Recompute diagonals $i+1, \ldots, N$ using Burge
    with the new vertex partitions~$\mu^{(i,\cdot)}$.  This
    produces a new ALCD~$\mathfrak{d}'$.
\end{enumerate}
Now apply $\sigma_i^Q$ again to $(S, \mathfrak{d}')$:
\begin{enumerate}
  \item[(a$'$)] Switch diagonal~$i$ back from RSK to Burge
    (since $\sigma_i^Q$ toggles the rule at diagonal~$i$).
    The vertex partitions at diagonal~$i$ are restored to
    the originals~$\lambda^{(i,\cdot)}$.
  \item[(b$'$)] Recompute diagonals $i+1, \ldots, N$ using
    Burge with the original vertex
    partitions~$\lambda^{(i,\cdot)}$.  This reproduces the
    original ALCD~$\mathfrak{d}$.
\end{enumerate}
The key point: step~(a$'$) restores the vertex partitions at
diagonal~$i$ because the local rule is a bijection, and
switching Burge $\to$ RSK $\to$ Burge at the same diagonal
with the same inputs below diagonal~$i$ (which are unchanged)
gives back the original output.  Step~(b$'$) then recomputes
the rest deterministically, recovering~$\mathfrak{d}$.
\end{proof}

\begin{remark}[The bijective property is essential]
\label{rem:bijective}
The proof uses the fact that both Burge and RSK are bijections
with the same type signature.  Switching from Burge to RSK at
diagonal~$i$ produces a well-defined output, and switching
back recovers the original.  This is the growth-diagram
analogue of the trivial fact that a toggle (Burge $\to$ RSK
$\to$ Burge) is the identity.

The propagation to subsequent diagonals is
\emph{automatically} self-inverse because the recomputation
step is deterministic: given the vertex partitions at
diagonal~$i$, the rest of the ALCD is uniquely determined
by the local rule.
\end{remark}


%% ============================================================
\section{The $Q$-Bender--Knuth involution}
\label{sec:Q-BK}
%% ============================================================

\begin{definition}[$Q$-Bender--Knuth]
\label{def:Q-BK}
For $1 \leq i \leq n - |S|$, the \emph{$Q$-Bender--Knuth
involution} is:
\[
  \sigma_i^Q(S, \mathfrak{d})
  = (S,\, \sigma_i^Q(\mathfrak{d})).
\]
\end{definition}

\begin{proposition}[Properties]
\label{prop:Q-BK-properties}
\mbox{}
\begin{enumerate}
  \item $(\sigma_i^Q)^2 = \mathrm{id}$
    \textup{(Theorem~\ref{thm:Q-sigma-involution})}.
  \item $\sigma_i^Q$ preserves the ground-set level, weight,
    and sign.
  \item $\sigma_i^Q \circ \sigma_j^Q
    = \sigma_j^Q \circ \sigma_i^Q$ for $|i - j| \geq 2$.
\end{enumerate}
\end{proposition}

\begin{proof}
(1)~Proved above.
(2)~The subset~$S$ is unchanged (sign and weight from~$S$
preserved), and the total ALCD weight is preserved by the
local-rule weight condition.
(3)~When $|i - j| \geq 2$, switching diagonal~$i$ changes
vertex partitions starting at diagonal~$i$, and switching
diagonal~$j$ changes vertex partitions starting at
diagonal~$j$.  Since the intervening diagonals are recomputed
using the same (original) rule, the two switches don't
interact.
\end{proof}


\subsection{Differences from the classical Bender--Knuth}

\begin{center}
\begin{tabular}{l|l}
\textbf{Classical $\sigma_i$} &
  \textbf{$Q$-enriched $\sigma_i^Q$} \\ \hline
Changes face labels at diagonal~$i$ only &
  Changes face labels at diagonals $\geq i$ \\
Free/paired decomposition determines everything &
  Propagation modifies downstream diagonals \\
$\sigma_i \sigma_{i+1} = \sigma_{i+1} \sigma_i$
  (if non-adjacent) &
  Same for $|i-j| \geq 2$ \\
Braid relation
  $\sigma_i \sigma_{i+1} \sigma_i
  = \sigma_{i+1} \sigma_i \sigma_{i+1}$ &
  \textbf{Open} (propagation may break it)
\end{tabular}
\end{center}

\begin{remark}[The braid relation]
\label{rem:braid}
In the classical case, the Bender--Knuth involutions
satisfy the braid relation (also called the Yang--Baxter
relation):
\[
  \sigma_i \sigma_{i+1} \sigma_i
  = \sigma_{i+1} \sigma_i \sigma_{i+1}
  \quad \text{for all } i.
\]
Together with $\sigma_i^2 = \mathrm{id}$ and the commutativity
for $|i-j| \geq 2$, this gives a symmetric group action
$S_N \to \mathrm{Aut}(\SSYT(\lambda))$.

For the $Q$-enriched $\sigma_i^Q$, the braid relation may fail
because the propagation from diagonal~$i$ to diagonal~$i+1$
creates a coupling that doesn't exist in the classical case.
The $Q$-weights at diagonal~$i+1$ depend on the strip sizes,
which depend on the vertex partitions at diagonal~$i$, which
were modified by $\sigma_i^Q$.

If the braid relation does hold, it would give a
\emph{$Q$-symmetric group action} on the ground set, with
far-reaching consequences for the representation theory of
the $Q$-kernel.
\end{remark}


%% ============================================================
\section{Decomposition of the $Q$-Sch\"utzenberger}
\label{sec:decomposition}
%% ============================================================

\begin{proposition}
\label{prop:psi-decomp}
The $Q$-Sch\"utzenberger $\psi_Q$ (switching all diagonals
from Burge to RSK, with recomputation) decomposes as:
\begin{equation}\label{eq:psi-Q-decomp}
  \psi_Q = \sigma_1^Q \circ \sigma_2^Q \circ \cdots
    \circ \sigma_N^Q,
\end{equation}
where $N = n - |S|$ is the number of ALCD diagonals and the
composition proceeds outside-in (switch diagonal~$N$ first).
\end{proposition}

\begin{theorem}[$\psi_Q^2 = \mathrm{id}$]
\label{thm:psi-Q-involution}
The $Q$-Sch\"utzenberger satisfies $\psi_Q^2 = \mathrm{id}$.
\end{theorem}

\begin{proof}
$\psi_Q$ recomputes the entire growth diagram using RSK
instead of Burge.  Applying $\psi_Q$ again recomputes using
Burge instead of RSK, restoring the original.  This is the
global toggle argument: the rule at each diagonal is switched
twice (Burge $\to$ RSK $\to$ Burge), which is the identity.

Note that this proof does \emph{not} require the
decomposition~\eqref{eq:psi-Q-decomp} or the braid relation.
It uses only the fact that $\psi_Q$ is a global
Burge$\leftrightarrow$RSK toggle.
\end{proof}

\begin{remark}[Two levels of involution]
\label{rem:two-levels}
We have involution properties at two levels:
\begin{enumerate}
  \item \textbf{Local}: $(\sigma_i^Q)^2 = \mathrm{id}$
    (switching one diagonal twice is the identity).
  \item \textbf{Global}: $\psi_Q^2 = \mathrm{id}$ (switching
    all diagonals twice is the identity).
\end{enumerate}
The global involution follows from the toggle structure without
needing the braid relation.  The braid relation would give
additional algebraic structure (a symmetric group action) but
is not needed for the $Q$-Sch\"utzenberger identity.
\end{remark}


%% ============================================================
\section{The $Q$-row switch on pairs}
\label{sec:Q-row-pairs}
%% ============================================================

For pairs $\mathcal{P}_{a,b,k} = \mathcal{G}_{a+k,c}
\times \mathcal{G}_{b+k,c}$, the $Q$-row switch acts on
each component:

\begin{definition}
\label{def:Q-row-pairs}
$\sigma_{i,j}^Q\bigl((S, \mathfrak{d}),\,
(S', \mathfrak{d}')\bigr)
= \bigl(\sigma_i^Q(S, \mathfrak{d}),\,
\sigma_j^Q(S', \mathfrak{d}')\bigr)$.
\end{definition}

The pair-level $Q$-Sch\"utzenberger
$\psi_Q = \mathfrak{R} \circ \mathfrak{B}^{-1}$
from~\cite{Q_SCHUTZ} corresponds to applying the full
Burge$\to$RSK switch on the first component and the full
RSK$\to$Burge switch on the second component.  By
Theorem~\ref{thm:psi-Q-involution} applied to each
component independently, $\psi_Q^2 = \mathrm{id}$ on pairs.

This resolves Conjecture~3.3 of~\cite{Q_SCHUTZ}
affirmatively: the $Q$-Sch\"utzenberger is indeed an
involution.


%% ============================================================
\section{Summary and open questions}
\label{sec:summary}
%% ============================================================

\begin{center}
\begin{tabular}{l|l|l}
\textbf{Classical} & \textbf{$Q$-enriched} & \textbf{Status}
  \\ \hline
$\sigma_i^2 = \mathrm{id}$ &
  $(\sigma_i^Q)^2 = \mathrm{id}$ &
  \checkmark\; proved \\
$\sigma_i$ local (one diagonal) &
  $\sigma_i^Q$ propagates downstream &
  structural difference \\
$\sigma_i \sigma_j = \sigma_j \sigma_i$, $|i-j| \geq 2$ &
  same &
  \checkmark\; proved \\
braid: $\sigma_i \sigma_{i+1} \sigma_i
  = \sigma_{i+1} \sigma_i \sigma_{i+1}$ &
  $\sigma_i^Q \sigma_{i+1}^Q \sigma_i^Q
  \stackrel{?}{=} \sigma_{i+1}^Q \sigma_i^Q
  \sigma_{i+1}^Q$ &
  \textbf{open} \\
$\psi = \prod \sigma_i$, $\psi^2 = \mathrm{id}$ &
  $\psi_Q = \prod \sigma_i^Q$, $\psi_Q^2 = \mathrm{id}$ &
  \checkmark\; proved \\
$\langle \sigma_1, \ldots, \sigma_N \rangle \cong S_{N+1}$ &
  $\langle \sigma_1^Q, \ldots, \sigma_N^Q \rangle
  \stackrel{?}{\cong} S_{N+1}$ &
  needs braid
\end{tabular}
\end{center}

\medskip
\noindent\textbf{Open questions.}
\begin{enumerate}
  \item \textbf{Braid relation.}\;
    Does $\sigma_i^Q \sigma_{i+1}^Q \sigma_i^Q
    = \sigma_{i+1}^Q \sigma_i^Q \sigma_{i+1}^Q$?
    If so, the $Q$-Bender--Knuth involutions generate a
    symmetric group action on the ground set, giving a
    $Q$-analogue of Corollary~\ref{cor:schur-symmetric}
    (a symmetry property of a ``$Q$-Schur function'').

  \item \textbf{Free/paired decomposition.}\;
    Is there a $Q$-analogue of the free/paired decomposition
    (Definition~\ref{def:free-paired}) that describes which
    face labels change under $\sigma_i^Q$ and which are
    fixed?  The propagation means that ``fixed'' labels may
    change at downstream diagonals even though they are
    ``paired'' at diagonal~$i$.

  \item \textbf{Computational verification.}\;
    For small $d$ and $n$, compute $\sigma_i^Q$ on explicit
    ground-set elements, verify $(\sigma_i^Q)^2 = \mathrm{id}$
    and test the braid relation.

  \item \textbf{Crystal-base connection.}\;
    The classical Bender--Knuth involutions correspond to
    Kashiwara's crystal operators.  Do the $\sigma_i^Q$
    correspond to crystal operators for an affine Lie algebra
    $\widehat{\mathfrak{sl}}_3$ (matching the rank-$3$
    cylindric partition structure)?
\end{enumerate}


\begin{thebibliography}{99}
\bibitem{ADJOINT}
  R.~Langer,
  \emph{The $Q$-Operator and its Adjoint: From Level~3
  Operator Traces to Cascade Commutators},
  companion note, 2026.

\bibitem{BenderKnuth}
  E.~A.~Bender and D.~E.~Knuth,
  \emph{Enumeration of plane partitions},
  J.~Combin.\ Theory Ser.~A \textbf{13} (1972), 40--54.

\bibitem{LOCAL_RULE}
  R.~Langer,
  \emph{The Generalized Local Rule and the Burge--RSK
  Involution},
  companion note, 2026.

\bibitem{Q_SCHUTZ}
  R.~Langer,
  \emph{The $Q$-Sch\"utzenberger Identity},
  companion note, 2026.

\bibitem{UNDERLYING_SET}
  R.~Langer,
  \emph{The Underlying Set: Cylindric Plane Partitions, ALCDs,
  and the Ground Set for the Involution},
  companion note.

\bibitem{langer-thesis}
  R.~Langer,
  \emph{Cylindric Partitions and Bijective Proofs},
  PhD thesis, University of Melbourne.

\bibitem{vanLeeuwen}
  M.~A.~A.~van~Leeuwen,
  \emph{The Robinson--Schensted and Sch\"utzenberger
  algorithms, an elementary approach},
  Electron.~J.~Combin.\ \textbf{3} (1996), R15.
\end{thebibliography}

\end{document}
